\documentclass[conference]{IEEEtran}
\usepackage{cite}
\usepackage{amsmath,amssymb,amsfonts}
\usepackage{graphicx}
\usepackage{textcomp}
\usepackage{xcolor}
\usepackage{booktabs}
\usepackage{listings}
\usepackage{hyperref}
\usepackage{cleveref}

\begin{document}

\title{Skill-Based Agentic Coding for Mathematical Software:\\
A Case Study in NP-Hard Problem Reductions}

\author{...}  % placeholder

\maketitle

\begin{abstract}
AI coding agents achieve 70--80\% on single-issue benchmarks like SWE-Bench Verified, but their success rate drops below 25\% on long-horizon software evolution tasks that demand sustained mathematical reasoning across many files. We address this gap by decomposing agentic coding into two complementary roles: human-creative work (designing reduction proofs, choosing algorithms, writing specifications) and agent-managed execution (scaffolding, testing, verification, and integration). Our method centers on a library of 13 reusable skills---from issue triage through implementation to multi-layered review---orchestrated by a coding agent within a Rust library for NP-hard problem reductions. A 7-layer verification stack (type checking, unit tests, brute-force cross-validation, closed-loop reduction tests, integration tests, coverage enforcement, and CI/CD) catches errors at increasing levels of abstraction. Applying this methodology over six months produced 24 problem types, 40 reduction rule implementations, and 52 edges in a typed reduction graph, all with $>$95\% test coverage. We contribute the skill-based decomposition methodology, the verification stack design, and the open-source artifact as a benchmark for agentic mathematical software engineering.
\end{abstract}

\section{Introduction}\label{sec:intro}

AI coding agents have made remarkable progress on isolated software engineering tasks.
On SWE-Bench Verified, which evaluates agents on single-issue bug fixes drawn from popular open-source repositories, the best systems now resolve 70--80\% of issues end-to-end~\cite{Xia2025LiveSWEagent}.
This has fueled optimism that fully autonomous software engineering is within reach.
Yet benchmarks designed to probe longer-horizon capabilities tell a starkly different story.
SWE-EVO, which requires agents to implement multi-step modifications spanning an average of 21~files, reports resolution rates around 21\% even for frontier models~\cite{Thai2025SWEEVO}.
SWE-Bench Pro, targeting enterprise-level tasks that may require hours to days of human effort, similarly finds that agents struggle with sustained multi-file reasoning~\cite{Deng2025SWEBenchPro}.
The common response to this capability gap has been to push for more powerful agents---larger models, better tool use, self-evolving scaffolds.
We argue that this framing overlooks the more fundamental bottleneck: not raw agent capability, but how work is decomposed and distributed between humans and agents.

Our thesis is that the key to effective agentic coding lies in \emph{task decomposition}: splitting the creative and judgment-intensive aspects of software development (which humans do well) from the management and mechanical aspects (which agents can handle reliably).
When a task is sufficiently well-specified and bounded in scope, even current agents execute it with high fidelity.
The challenge is not making agents smarter, but structuring the work so that each unit falls within the agent's reliable operating range.
This perspective shifts attention from agent architecture to \emph{skill design}---the craft of encoding domain knowledge into reusable, agent-executable task specifications.

This decomposition is especially critical for mathematical and scientific software, where the ``review is harder than generation'' problem is acute~\cite{Roychoudhury2025AgenticAI}.
An agent can generate a plausible-looking reduction from Boolean satisfiability to graph coloring, but verifying that the reduction preserves solution structure requires mathematical reasoning that current agents cannot reliably perform in isolation.
Roychoudhury~\cite{Roychoudhury2025AgenticAI} identifies specification inference---deciphering and clarifying developer intent---as the central difficulty in agentic software workflows, and argues that trustworthy deployment requires AI-based verification and validation of AI-generated code.
Our multi-layered verification stack addresses precisely this challenge: rather than attempting end-to-end formal verification (which remains largely unsolved for complex mathematical code~\cite{Thakur2025CLEVER, Bursuc2025VeriCoding}), we compose multiple lightweight verification mechanisms that collectively catch errors across different abstraction levels.

We instantiate this methodology in the domain of NP-hard problem reductions, implemented as an open-source Rust library.
The library manages a typed reduction graph connecting 24 problem types through 40 hand-coded reduction rule implementations and 52 total directed edges (including 12 edges inferred from a type-parameter subtype lattice).
The domain serves as a Goldilocks testbed for studying agentic coding: each reduction is self-contained (50--200 lines of code), requires non-trivial mathematical reasoning about the mapping between problem structures, yet admits a fully automatable correctness criterion---reduce an instance, solve the target problem by brute force, extract the solution back, and verify it against the source.
This combination of mathematical depth with mechanical verifiability makes it possible to study how agents perform on tasks that are individually tractable but collectively demand sustained engineering discipline.

Our approach distributes work across three roles with distinct responsibilities.
\textbf{Contributors} perform the creative work of identifying which problems and reductions are worth implementing: they open issues proposing new nodes or edges in the reduction graph, drawing on domain knowledge to spot gaps, recognize useful connections, and assess mathematical non-triviality.
\textbf{The maintainer} curates the project board and writes skills---markdown scripts that decompose complex tasks (such as ``implement a new reduction rule'') into sequences of agent-manageable subtasks.
The maintainer encodes domain knowledge, quality standards, and project conventions into these skills, effectively programming the agent's workflow rather than its output.
\textbf{Agents} serve in a dual capacity: as \emph{managers}, they pick cards from the project board, dispatch sub-agents for parallel review, and orchestrate a two-stage pipeline from issue to merged pull request; as \emph{executors}, they implement code, write tests, generate documentation, and fix CI failures.
The key insight is that the two human roles contribute \emph{judgment}---which reductions matter, what quality bar to enforce---while the agent handles \emph{volume}---executing the mechanical steps reliably and repeatedly.
Industry data supports this division: developers now use AI in 60\% of their work while maintaining active oversight on 80--100\% of delegated tasks~\cite{Anthropic2026AgenticCoding}.

We organize our agent's capabilities into a library of 13~skills spanning five functional categories: orchestration (pipeline management and issue dispatch), implementation (adding models and reduction rules), quality gates (issue checking, redundancy analysis, multi-agent review, and CI repair), documentation (generating formal problem definitions and reduction theorems with proof sketches), and release management.
A two-stage card-based pipeline automates the progression from issue to merged code: the first stage picks a ``Ready'' issue, implements it in an isolated git worktree, and produces a pull request; the second stage addresses review comments, fixes CI failures, and prepares the PR for human merge.
The human maintainer touches only two transitions in this pipeline---moving an issue from Backlog to Ready (the strategic decision of \emph{what} to work on) and merging the final pull request (the quality gate of \emph{whether} the work meets standards).
Everything in between is agent-managed.

Correctness assurance comes from a seven-layer verification stack that catches errors at increasing levels of abstraction.
The stack ranges from compile-time type checking (Layer~1), through unit tests and brute-force cross-validation (Layers~2--3), to overhead formula validation against actual reduction sizes (Layer~4), materialized test fixtures that prevent agents from silently changing expected values (Layer~5), parallel agentic review with fresh context windows (Layer~6), and finally, documentation entries that require human-readable proof sketches for each reduction (Layer~7).
No single layer is sufficient---the type system catches API misuse but not logical errors; closed-loop tests verify functional correctness but not overhead formulas; documentation catches proof-level mistakes that no automated test can detect.
The layers are designed to be complementary, and the skill system ensures that agents invoke all relevant layers as part of every implementation task.

Our contributions are as follows:
\begin{itemize}
  \item A \textbf{skill-based methodology} for decomposing mathematical coding tasks into agent-manageable steps, with a concrete skill library and a card-based orchestration pipeline that separates human judgment from agent execution.
  \item A \textbf{multi-layered verification stack} comprising seven complementary mechanisms---from type-level enforcement through materialized fixtures to agentic review and documentation-as-verification---that collectively ensure correctness of agent-generated mathematical code.
  \item A \textbf{verified open-source artifact}: a Rust library implementing 24~NP-hard problem types, 40~reduction rules, and 52~graph edges, all with ${>}95\%$ test coverage, serving as both a practical tool for reduction-based problem solving and a benchmark for evaluating agentic mathematical software engineering.
\end{itemize}

The remainder of this paper is organized as follows.
\Cref{sec:domain} motivates our choice of NP-hard reductions as a study domain, arguing that it occupies a sweet spot between mathematical complexity and mechanical verifiability.
\Cref{sec:architecture} describes the Rust library's type-driven architecture that makes agent-generated code verifiable by construction.
\Cref{sec:skills} presents the skill-based task decomposition, the three-role model, and the card-based orchestration pipeline.
\Cref{sec:verification} details the seven-layer verification stack.
\Cref{sec:evaluation} evaluates the methodology through an ablation study, git history mining, and detailed case studies.
\Cref{sec:related} surveys related work on AI coding agents, AI-assisted discovery of reductions, and formal verification of generated code.
\Cref{sec:conclusion} discusses generalizability, limitations, and future directions.

\section{Why Reductions? The Goldilocks Domain}\label{sec:domain}

\begin{figure*}[t]
  \centering
  \includegraphics[width=\textwidth]{figures/reduction-graph.pdf}
  \caption{The reduction graph: 24 problem types connected by 52 directed edges (35 implemented reductions + 17 inferred variant edges). Hub nodes MIS, QUBO and ILP are highlighted. Nodes in dashed boxes have models but no cross-type reductions yet.}
  \label{fig:reduction-graph}
\end{figure*}

\section{System Architecture}\label{sec:architecture}

\begin{figure}[t]
  \centering
  \includegraphics[width=\columnwidth]{figures/architecture.pdf}
  \caption{System architecture: the trait hierarchy and compile-time validation enforce round-trip testing capability by construction.}
  \label{fig:architecture}
\end{figure}

\section{Skill-Based Task Decomposition}\label{sec:skills}

\begin{figure}[t]
  \centering
  \includegraphics[width=\columnwidth]{figures/pipeline.pdf}
  \caption{Two-stage card-based pipeline. Human decisions (orange) are limited to Backlog$\to$Ready and In Review$\to$Done. Agent manages everything in between.}
  \label{fig:pipeline}
\end{figure}

\section{Multi-Layered Verification}\label{sec:verification}

\begin{figure}[t]
  \centering
  \includegraphics[width=\columnwidth]{figures/verification-pyramid.pdf}
  \caption{Seven-layer verification stack. Lower layers (blue) are fully automated; upper layers (gold) involve human-readable arguments.}
  \label{fig:verification}
\end{figure}

\section{Evaluation}\label{sec:evaluation}

\section{Related Work}\label{sec:related}

\section{Discussion \& Conclusion}\label{sec:conclusion}

\bibliographystyle{IEEEtran}
\bibliography{references}

\end{document}
